
\begin{center}
 {\large {\textbf 
5º Abóbora Contest}}

%Sub-Regional Brasil do ACM ICPC

{\large 31 de Outubro de 2024}

{\large  Caderno de Problemas}

{\large \textbf{Informações Gerais}}
\end{center}

Este caderno contém $8$ problemas; as páginas estão numeradas de $1$ a $10$,  contando esta página de rosto. Verifique se o caderno está completo.

\begin{enumerate}
    \item \textbf{Sobre os nomes dos programas}
    \begin{enumerate}
        \item Para soluções em C/C++ e Python 3.10, o nome do arquivo-fonte não é significativo, pode ser qualquer nome desde que tenha as extensões .c, .cpp, .java e .py.
        \item Para linguagem Java, sua solução deve ser chamado de codigo\_do\_problema.java, em que codigo\_do\_problema é a letra maiúscula que identifica o problema. Lembre que em Java o nome da classe principal deve ser igual ao nome do arquivo.
    \end{enumerate}
    \item \textbf{Sobre a entrada}
    \begin{enumerate}
        \item A entrada de seu programa deve ser lida da entrada padrão.
        \item A entrada é composta de um único caso de teste, descrito em um número de linhas que depende do problema.
        \item Quando uma linha da entrada contém vários valores, estes são separados por um único espaço em branco; a entrada não contém nenhum outro espaço em branco.
        \item Cada linha, incluindo a última, contém exatamente um caractere final-de-linha.
        \item O final da entrada coincide com o final do arquivo.
    \end{enumerate}
    \item \textbf{Sobre a saída}
    \begin{enumerate}
        \item A saída de seu programa deve ser escrita na saída padrão.
        \item Quando uma linha da saída contém vários valores, estes devem ser separados por um único espaço em branco; a saída não deve conter nenhum outro espaço em branco.
        \item Cada linha, incluindo a última, deve conter exatamente um caractere final-de-linha.
    \end{enumerate}
\end{enumerate}

\newpage